%% Response to the reviewers, revision 2 of CAOR-D-26-01885.
%% Build (after compiling ../main.tex, whose .aux supplies the cross
%% references): pdflatex response_letter (twice)
\documentclass[11pt]{article}
\usepackage[a4paper,margin=2.4cm]{geometry}
\usepackage{amsmath,amssymb}
\usepackage{xcolor}
\usepackage{parskip}
\usepackage{enumitem}
\usepackage{xr}
\usepackage[colorlinks=true,allcolors=blue!50!black]{hyperref}
\externaldocument[m-]{../main}
% generated by make_tables.py — do not hand-edit
\newcommand{\batonGrandBest}{53.7\%}
\newcommand{\budCityDp}{12.2\%}
\newcommand{\budCityDpNLarge}{11.9\%}
\newcommand{\budCityDpNMid}{7.0\%}
\newcommand{\budCityDpNSmall}{-0.1\%}
\newcommand{\budCityDpThree}{12.4\%}
\newcommand{\budCityLarge}{12.9\%}
\newcommand{\budCityMid}{10.4\%}
\newcommand{\budCitySmall}{3.9\%}
\newcommand{\budCityThrLarge}{12.0\%}
\newcommand{\budCityThrMid}{9.9\%}
\newcommand{\budCityThrSmall}{0.5\%}
\newcommand{\budDethDetDp}{26.4\%}
\newcommand{\budDethDetDpNLarge}{26.3\%}
\newcommand{\budDethDetDpNMid}{25.7\%}
\newcommand{\budDethDetDpNSmall}{22.7\%}
\newcommand{\budDethDetDpThree}{29.5\%}
\newcommand{\budDethDetLarge}{28.3\%}
\newcommand{\budDethDetMid}{27.9\%}
\newcommand{\budDethDetSmall}{25.8\%}
\newcommand{\budDethDetThrLarge}{25.0\%}
\newcommand{\budDethDetThrMid}{24.8\%}
\newcommand{\budDethDetThrSmall}{23.0\%}
\newcommand{\budDethSaaDp}{22.2\%}
\newcommand{\budDethSaaDpNLarge}{22.0\%}
\newcommand{\budDethSaaDpNMid}{12.1\%}
\newcommand{\budDethSaaDpNSmall}{-1.6\%}
\newcommand{\budDethSaaDpThree}{56.4\%}
\newcommand{\budDethSaaLarge}{56.3\%}
\newcommand{\budDethSaaMid}{53.6\%}
\newcommand{\budDethSaaSmall}{41.7\%}
\newcommand{\budDethSaaThrLarge}{19.6\%}
\newcommand{\budDethSaaThrMid}{16.5\%}
\newcommand{\budDethSaaThrSmall}{-2.4\%}
\newcommand{\budgetCross}{100}
\newcommand{\certBeatN}{17/40}
\newcommand{\certGap}{9.34\%}
\newcommand{\certGapHighs}{11.03\%}
\newcommand{\certGapMax}{12.55\%}
\newcommand{\cfGapHi}{1.3}
\newcommand{\cfGapLo}{0.1}
\newcommand{\cityActShare}{37\%}
\newcommand{\cityDetK}{8.5}
\newcommand{\cityDetReact}{37.2}
\newcommand{\citySaaK}{12.1}
\newcommand{\citySaaPem}{0.05\%}
\newcommand{\citySaaReact}{0.04}
\newcommand{\costCuiBaton}{11.5}
\newcommand{\costDetBaton}{135.0}
\newcommand{\costGounarisBaton}{10.2}
\newcommand{\costMDROBaton}{2.2}
\newcommand{\costSAABaton}{8.3}
\newcommand{\costWDROBaton}{1.9}
\newcommand{\csBaseBaton}{52.7}
\newcommand{\csBaseReact}{78.4}
\newcommand{\csEmgHighBaton}{57.7}
\newcommand{\csEmgHighReact}{97.3}
\newcommand{\csEmgHighSv}{55.2\%}
\newcommand{\csEmgLowBaton}{48.3}
\newcommand{\csEmgLowReact}{64.2}
\newcommand{\csEmgLowSv}{40.1\%}
\newcommand{\csSbHighBaton}{60.9}
\newcommand{\csSbHighReact}{78.4}
\newcommand{\csSbHighSv}{42.6\%}
\newcommand{\csSbLowBaton}{45.8}
\newcommand{\csSbLowReact}{78.4}
\newcommand{\csSbLowSv}{52.7\%}
\newcommand{\ctdBaton}{51.9\%}
\newcommand{\ctdBatonFull}{69.1\%}
\newcommand{\ctdVOne}{0.0\%}
\newcommand{\dearSbBaton}{42.6\%}
\newcommand{\dearSbOrc}{24.8\%}
\newcommand{\dearSbThresh}{8.8\%}
\newcommand{\depCfAbove}{8}
\newcommand{\depDayfacBaton}{10.8\%}
\newcommand{\depDayfacDetBaton}{14.2\%}
\newcommand{\depDayfacDetCf}{14.2\%}
\newcommand{\depDayfacDetComp}{two-lever}
\newcommand{\depDayfacDetCompSv}{11.8\%}
\newcommand{\depDayfacDetDelta}{+2.3}
\newcommand{\depDayfacDetDeltaCI}{[+1.4, +3.1]}
\newcommand{\depDayfacDetDpThree}{14.2\%}
\newcommand{\depDayfacDetHo}{9.8\%}
\newcommand{\depDayfacDetRatio}{100\%}
\newcommand{\depDayfacDetReact}{172.2}
\newcommand{\depDayfacDetThr}{9.1\%}
\newcommand{\depDayfacDetThrK}{9.7\%}
\newcommand{\depDayfacSaaBaton}{7.4\%}
\newcommand{\depDayfacSaaCf}{8.0\%}
\newcommand{\depDayfacSaaComp}{restock}
\newcommand{\depDayfacSaaCompSv}{5.8\%}
\newcommand{\depDayfacSaaDelta}{+1.6}
\newcommand{\depDayfacSaaDeltaCI}{[-1.6, +4.9]}
\newcommand{\depDayfacSaaDpThree}{12.4\%}
\newcommand{\depDayfacSaaHo}{-6.9\%}
\newcommand{\depDayfacSaaRatio}{60\%}
\newcommand{\depDayfacSaaReact}{5.2}
\newcommand{\depDayfacSaaThr}{-15.6\%}
\newcommand{\depDayfacSaaThrK}{-12.7\%}
\newcommand{\depNClose}{4}
\newcommand{\depRhoNineBaton}{49.9\%}
\newcommand{\depRhoNineDetBaton}{35.4\%}
\newcommand{\depRhoNineDetCf}{37.6\%}
\newcommand{\depRhoNineDetComp}{DP$^3_N$}
\newcommand{\depRhoNineDetCompSv}{36.5\%}
\newcommand{\depRhoNineDetDelta}{-1.1}
\newcommand{\depRhoNineDetDeltaCI}{[-1.6, -0.7]}
\newcommand{\depRhoNineDetDpThree}{38.0\%}
\newcommand{\depRhoNineDetHo}{34.1\%}
\newcommand{\depRhoNineDetRatio}{93\%}
\newcommand{\depRhoNineDetReact}{177.5}
\newcommand{\depRhoNineDetThr}{32.3\%}
\newcommand{\depRhoNineDetThrK}{34.0\%}
\newcommand{\depRhoNineSaaBaton}{64.5\%}
\newcommand{\depRhoNineSaaCf}{65.9\%}
\newcommand{\depRhoNineSaaComp}{two-lever}
\newcommand{\depRhoNineSaaCompSv}{54.4\%}
\newcommand{\depRhoNineSaaDelta}{+10.2}
\newcommand{\depRhoNineSaaDeltaCI}{[+8.8, +11.6]}
\newcommand{\depRhoNineSaaDpThree}{67.7\%}
\newcommand{\depRhoNineSaaHo}{33.0\%}
\newcommand{\depRhoNineSaaRatio}{95\%}
\newcommand{\depRhoNineSaaReact}{25.9}
\newcommand{\depRhoNineSaaThr}{27.3\%}
\newcommand{\depRhoNineSaaThrK}{32.9\%}
\newcommand{\depRhoSixBaton}{40.8\%}
\newcommand{\depRhoSixDetBaton}{27.9\%}
\newcommand{\depRhoSixDetCf}{29.2\%}
\newcommand{\depRhoSixDetComp}{DP$^3_N$}
\newcommand{\depRhoSixDetCompSv}{27.8\%}
\newcommand{\depRhoSixDetDelta}{+0.0}
\newcommand{\depRhoSixDetDeltaCI}{[-0.4, +0.4]}
\newcommand{\depRhoSixDetDpThree}{29.5\%}
\newcommand{\depRhoSixDetHo}{26.1\%}
\newcommand{\depRhoSixDetRatio}{94\%}
\newcommand{\depRhoSixDetReact}{179.2}
\newcommand{\depRhoSixDetThr}{24.7\%}
\newcommand{\depRhoSixDetThrK}{26.0\%}
\newcommand{\depRhoSixSaaBaton}{53.7\%}
\newcommand{\depRhoSixSaaCf}{53.8\%}
\newcommand{\depRhoSixSaaComp}{two-lever}
\newcommand{\depRhoSixSaaCompSv}{43.9\%}
\newcommand{\depRhoSixSaaDelta}{+9.7}
\newcommand{\depRhoSixSaaDeltaCI}{[+8.2, +11.1]}
\newcommand{\depRhoSixSaaDpThree}{56.4\%}
\newcommand{\depRhoSixSaaHo}{20.1\%}
\newcommand{\depRhoSixSaaRatio}{95\%}
\newcommand{\depRhoSixSaaReact}{16.5}
\newcommand{\depRhoSixSaaThr}{16.3\%}
\newcommand{\depRhoSixSaaThrK}{19.8\%}
\newcommand{\depRhoThreeBaton}{29.1\%}
\newcommand{\depRhoThreeDetBaton}{21.9\%}
\newcommand{\depRhoThreeDetCf}{22.5\%}
\newcommand{\depRhoThreeDetComp}{DP$^3_N$}
\newcommand{\depRhoThreeDetCompSv}{20.7\%}
\newcommand{\depRhoThreeDetDelta}{+1.2}
\newcommand{\depRhoThreeDetDeltaCI}{[+0.9, +1.6]}
\newcommand{\depRhoThreeDetDpThree}{22.7\%}
\newcommand{\depRhoThreeDetHo}{19.3\%}
\newcommand{\depRhoThreeDetRatio}{97\%}
\newcommand{\depRhoThreeDetReact}{180.2}
\newcommand{\depRhoThreeDetThr}{18.0\%}
\newcommand{\depRhoThreeDetThrK}{19.1\%}
\newcommand{\depRhoThreeSaaBaton}{36.3\%}
\newcommand{\depRhoThreeSaaCf}{36.6\%}
\newcommand{\depRhoThreeSaaComp}{restock}
\newcommand{\depRhoThreeSaaCompSv}{29.8\%}
\newcommand{\depRhoThreeSaaDelta}{+6.5}
\newcommand{\depRhoThreeSaaDeltaCI}{[+4.7, +8.4]}
\newcommand{\depRhoThreeSaaDpThree}{40.9\%}
\newcommand{\depRhoThreeSaaHo}{7.9\%}
\newcommand{\depRhoThreeSaaRatio}{89\%}
\newcommand{\depRhoThreeSaaReact}{8.0}
\newcommand{\depRhoThreeSaaThr}{3.1\%}
\newcommand{\depRhoThreeSaaThrK}{6.5\%}
\newcommand{\depRhoZeroBaton}{7.7\%}
\newcommand{\depRhoZeroDetBaton}{15.5\%}
\newcommand{\depRhoZeroDetCf}{15.4\%}
\newcommand{\depRhoZeroDetComp}{DP$^3_N$}
\newcommand{\depRhoZeroDetCompSv}{13.5\%}
\newcommand{\depRhoZeroDetDelta}{+2.0}
\newcommand{\depRhoZeroDetDeltaCI}{[+1.5, +2.1]}
\newcommand{\depRhoZeroDetDpThree}{15.5\%}
\newcommand{\depRhoZeroDetHo}{10.6\%}
\newcommand{\depRhoZeroDetRatio}{100\%}
\newcommand{\depRhoZeroDetReact}{176.4}
\newcommand{\depRhoZeroDetThr}{9.9\%}
\newcommand{\depRhoZeroDetThrK}{10.5\%}
\newcommand{\depRhoZeroSaaBaton}{-0.1\%}
\newcommand{\depRhoZeroSaaCf}{0.2\%}
\newcommand{\depRhoZeroSaaComp}{DP$_N$}
\newcommand{\depRhoZeroSaaCompSv}{0.0\%}
\newcommand{\depRhoZeroSaaDelta}{-0.1}
\newcommand{\depRhoZeroSaaDeltaCI}{[-4.6, +2.2]}
\newcommand{\depRhoZeroSaaDpThree}{7.4\%}
\newcommand{\depRhoZeroSaaHo}{-8.0\%}
\newcommand{\depRhoZeroSaaRatio}{-2\%}
\newcommand{\depRhoZeroSaaReact}{2.8}
\newcommand{\depRhoZeroSaaThr}{-24.9\%}
\newcommand{\depRhoZeroSaaThrK}{-11.9\%}
\newcommand{\detActGain}{1.8}
\newcommand{\dtAllAware}{27.6\%}
\newcommand{\dtAllPooled}{27.6\%}
\newcommand{\dtAwarePooledCI}{[-0.7, -0.3]}
\newcommand{\dtAwarePooledDelta}{-0.5}
\newcommand{\dtAwarePooledDetCI}{[-0.1, +0.3]}
\newcommand{\dtAwarePooledDetDelta}{+0.1}
\newcommand{\dtAwarePooledSaaCI}{[-1.5, -0.8]}
\newcommand{\dtAwarePooledSaaDelta}{-1.1}
\newcommand{\dtNNormal}{803}
\newcommand{\dtNPromo}{197}
\newcommand{\dtPromoAware}{29.7\%}
\newcommand{\dtPromoPooled}{29.4\%}
\newcommand{\dtPromoStale}{28.3\%}
\newcommand{\dtStalePooledCI}{[-1.7, -1.0]}
\newcommand{\dtStalePooledDelta}{-1.3}
\newcommand{\dtStalePooledDetCI}{[-0.6, -0.2]}
\newcommand{\dtStalePooledDetDelta}{-0.4}
\newcommand{\dtStalePooledSaaCI}{[-2.9, -1.6]}
\newcommand{\dtStalePooledSaaDelta}{-2.2}
\newcommand{\exDetBaton}{12.6\%}
\newcommand{\exDetBatonHo}{9.0\%}
\newcommand{\exDetDpThree}{12.6\%}
\newcommand{\exDetHo}{9.6\%}
\newcommand{\exDetMcErr}{0.39\%}
\newcommand{\exDetShare}{94\%}
\newcommand{\exDetShareDp}{94\%}
\newcommand{\exDetShareHo}{93\%}
\newcommand{\exDetThree}{13.4\%}
\newcommand{\exSaaBaton}{5.9\%}
\newcommand{\exSaaBatonHo}{-2.3\%}
\newcommand{\exSaaDpThree}{8.6\%}
\newcommand{\exSaaHo}{0.4\%}
\newcommand{\exSaaMcErr}{0.22\%}
\newcommand{\exSaaShare}{59\%}
\newcommand{\exSaaShareDp}{86\%}
\newcommand{\exSaaShareHo}{-607\%}
\newcommand{\exSaaThree}{10.0\%}
\newcommand{\feeBatonMax}{29.0\%}
\newcommand{\feeBreakEven}{29}
\newcommand{\feeGapMax}{0.0}
\newcommand{\feeGapZero}{19.1}
\newcommand{\feeHoMax}{29.0\%}
\newcommand{\feeMax}{30}
\newcommand{\frCityBiasIn}{-1.23\%}
\newcommand{\frCityBiasOut}{-1.09\%}
\newcommand{\frCityCf}{9.2\%}
\newcommand{\frCityDeltaF}{+0.2}
\newcommand{\frCityDeltaFCI}{[+0.1, +0.5]}
\newcommand{\frCityDpThree}{10.5\%}
\newcommand{\frCityDsEff}{-1.09\%}
\newcommand{\frCityDsEffCI}{[-1.7, -0.5]}
\newcommand{\frCityExact}{9.7\%}
\newcommand{\frCityFout}{5.4\%}
\newcommand{\frCityFstar}{5.7\%}
\newcommand{\frCityHo}{9.0\%}
\newcommand{\frCityInEff}{-0.13\%}
\newcommand{\frCityInEffCI}{[-0.5, +0.2]}
\newcommand{\frCityRsBaton}{0.8\%}
\newcommand{\frCityRsExact}{2.0\%}
\newcommand{\frCityThree}{5.5\%}
\newcommand{\frDethBiasIn}{+0.00\%}
\newcommand{\frDethBiasOut}{+0.01\%}
\newcommand{\frDethCf}{52.2\%}
\newcommand{\frDethDeltaF}{+0.0}
\newcommand{\frDethDeltaFCI}{[-0.04, +0.09]}
\newcommand{\frDethDpThree}{54.6\%}
\newcommand{\frDethDsEff}{+0.01\%}
\newcommand{\frDethDsEffCI}{[+0.00, +0.02]}
\newcommand{\frDethExact}{54.6\%}
\newcommand{\frDethFout}{51.4\%}
\newcommand{\frDethFstar}{51.4\%}
\newcommand{\frDethHo}{18.9\%}
\newcommand{\frDethInEff}{-0.01\%}
\newcommand{\frDethInEffCI}{[-0.05, +0.04]}
\newcommand{\frDethRsBaton}{4.1\%}
\newcommand{\frDethRsExact}{4.1\%}
\newcommand{\frDethThree}{51.4\%}
\newcommand{\frSalhiBiasIn}{+0.49\%}
\newcommand{\frSalhiBiasOut}{+0.09\%}
\newcommand{\frSalhiCf}{53.5\%}
\newcommand{\frSalhiDeltaF}{-0.3}
\newcommand{\frSalhiDeltaFCI}{[-0.77, -0.02]}
\newcommand{\frSalhiDpThree}{53.9\%}
\newcommand{\frSalhiDsEff}{+0.09\%}
\newcommand{\frSalhiDsEffCI}{[-0.5, +0.6]}
\newcommand{\frSalhiExact}{58.2\%}
\newcommand{\frSalhiFout}{52.6\%}
\newcommand{\frSalhiFstar}{52.2\%}
\newcommand{\frSalhiHo}{41.6\%}
\newcommand{\frSalhiInEff}{+0.40\%}
\newcommand{\frSalhiInEffCI}{[+0.2, +0.6]}
\newcommand{\frSalhiRsBaton}{20.1\%}
\newcommand{\frSalhiRsExact}{28.3\%}
\newcommand{\frSalhiThree}{52.6\%}
\newcommand{\gridCoarseHoMax}{2.93\%}
\newcommand{\gridCoarseThreeMax}{2.91\%}
\newcommand{\gridCoarseThreeMedian}{0.4\%}
\newcommand{\gridDetCoarseHo}{9.6\%}
\newcommand{\gridDetCoarseThree}{13.4\%}
\newcommand{\gridDetFineHo}{9.6\%}
\newcommand{\gridDetFineThree}{13.4\%}
\newcommand{\gridDetWideHo}{9.6\%}
\newcommand{\gridDetWideThree}{13.4\%}
\newcommand{\gridMidHoMax}{1.23\%}
\newcommand{\gridMidThreeMax}{1.20\%}
\newcommand{\gridPooledMaxDiff}{0.01}
\newcommand{\gridRoutes}{627}
\newcommand{\gridSaaCoarseHo}{0.4\%}
\newcommand{\gridSaaCoarseThree}{10.0\%}
\newcommand{\gridSaaFineHo}{0.4\%}
\newcommand{\gridSaaFineThree}{10.0\%}
\newcommand{\gridSaaWideHo}{0.4\%}
\newcommand{\gridSaaWideThree}{10.0\%}
\newcommand{\gridWideHoMax}{1.95\%}
\newcommand{\gridWideThreeMax}{1.95\%}
\newcommand{\lgCityBaton}{10.7\%}
\newcommand{\lgCityComp}{thr.-$k$}
\newcommand{\lgCityCompSv}{10.2\%}
\newcommand{\lgCityDelta}{+0.5}
\newcommand{\lgCityDeltaCI}{[+0.1, +0.7]}
\newcommand{\lgCityRatio}{86\%}
\newcommand{\lgCityRatioOrcThree}{28\%}
\newcommand{\lgCityUBaton}{11.2\%}
\newcommand{\lgCityUComp}{roll.-$\theta$}
\newcommand{\lgCityUCompSv}{11.1\%}
\newcommand{\lgCityUDelta}{+0.0}
\newcommand{\lgCityUDeltaCI}{[-0.8, +0.6]}
\newcommand{\lgCityURatio}{89\%}
\newcommand{\lgCityURatioOrcThree}{32\%}
\newcommand{\lgSalhiBaton}{53.3\%}
\newcommand{\lgSalhiComp}{DP$^3_N$}
\newcommand{\lgSalhiCompSv}{52.6\%}
\newcommand{\lgSalhiDelta}{+0.7}
\newcommand{\lgSalhiDeltaCI}{[+0.2, +1.2]}
\newcommand{\lgSalhiRatio}{98\%}
\newcommand{\lgSalhiRatioOrcThree}{85\%}
\newcommand{\lgZpFiftyBaton}{3.5\%}
\newcommand{\lgZpFiftyComp}{DP$_N$}
\newcommand{\lgZpFiftyCompSv}{3.6\%}
\newcommand{\lgZpFiftyDelta}{-0.1}
\newcommand{\lgZpFiftyDeltaCI}{[-7.1, +2.4]}
\newcommand{\lgZpFiftyRatio}{50\%}
\newcommand{\lgZpFiftyRatioOrcThree}{9\%}
\newcommand{\lgZpTwentyFiveBaton}{5.0\%}
\newcommand{\lgZpTwentyFiveComp}{thr.-$k$}
\newcommand{\lgZpTwentyFiveCompSv}{4.5\%}
\newcommand{\lgZpTwentyFiveDelta}{+0.4}
\newcommand{\lgZpTwentyFiveDeltaCI}{[+0.05, +0.62]}
\newcommand{\lgZpTwentyFiveRatio}{75\%}
\newcommand{\lgZpTwentyFiveRatioOrcThree}{15\%}
\newcommand{\nBatonTop}{six}
\newcommand{\nCfAbove}{all}
\newcommand{\nCity}{19}
\newcommand{\nOrcBeat}{four}
\newcommand{\nRL}{49}
\newcommand{\nRegretRoutes}{735}
\newcommand{\nTimingRoutes}{102}
\newcommand{\pemCui}{11.6\%}
\newcommand{\pemDet}{61.7\%}
\newcommand{\pemGounaris}{11.3\%}
\newcommand{\pemMDRO}{2.7\%}
\newcommand{\pemSAA}{10.0\%}
\newcommand{\pemWDRO}{2.3\%}
\newcommand{\poolBatonMax}{7}
\newcommand{\poolBatonMed}{0}
\newcommand{\poolCityBeatLow}{0/19}
\newcommand{\poolCityBreakEven}{2.6}
\newcommand{\poolCityGainOne}{5.1\%}
\newcommand{\poolCityGainZero}{10.6\%}
\newcommand{\poolCityK}{8}
\newcommand{\poolCityNaiveOne}{36.4}
\newcommand{\poolCityNaiveZero}{41.5}
\newcommand{\poolCityPpu}{33.8}
\newcommand{\poolCityShadowOne}{34.5}
\newcommand{\poolCityShadowZero}{37.1}
\newcommand{\poolCitySstar}{0.0}
\newcommand{\poolCitySstarLow}{0.1}
\newcommand{\poolCityTotHigh}{37.1}
\newcommand{\poolCityTotLow}{37.0}
\newcommand{\poolCuiMax}{1}
\newcommand{\poolCuiMed}{0}
\newcommand{\poolDetMax}{7}
\newcommand{\poolDetMed}{3.5}
\newcommand{\poolDethBeatLow}{40/40}
\newcommand{\poolDethBreakEven}{19.9}
\newcommand{\poolDethGainOne}{4.2\%}
\newcommand{\poolDethGainZero}{5.0\%}
\newcommand{\poolDethK}{7}
\newcommand{\poolDethNaiveOne}{154.3}
\newcommand{\poolDethNaiveZero}{176.5}
\newcommand{\poolDethPpu}{135.0}
\newcommand{\poolDethSaaBeatLow}{21/40}
\newcommand{\poolDethSaaBreakEven}{0.1}
\newcommand{\poolDethSaaGainOne}{0.6\%}
\newcommand{\poolDethSaaGainZero}{0.3\%}
\newcommand{\poolDethSaaK}{10}
\newcommand{\poolDethSaaNaiveOne}{8.2}
\newcommand{\poolDethSaaNaiveZero}{8.3}
\newcommand{\poolDethSaaPpu}{8.3}
\newcommand{\poolDethSaaShadowOne}{8.1}
\newcommand{\poolDethSaaShadowZero}{8.3}
\newcommand{\poolDethSaaSstar}{0.0}
\newcommand{\poolDethSaaSstarLow}{0.0}
\newcommand{\poolDethSaaTotHigh}{8.3}
\newcommand{\poolDethSaaTotLow}{8.3}
\newcommand{\poolDethShadowOne}{147.8}
\newcommand{\poolDethShadowZero}{167.7}
\newcommand{\poolDethSstar}{1.1}
\newcommand{\poolDethSstarLow}{4.2}
\newcommand{\poolDethTotHigh}{162.5}
\newcommand{\poolDethTotLow}{119.0}
\newcommand{\poolGounarisMax}{1}
\newcommand{\poolGounarisMed}{0}
\newcommand{\poolHoMax}{8}
\newcommand{\poolHoMed}{1}
\newcommand{\poolMDROMax}{0}
\newcommand{\poolMDROMed}{0}
\newcommand{\poolMetroBeatLow}{4/4}
\newcommand{\poolMetroBreakEven}{6.6}
\newcommand{\poolMetroGainMax}{9.4\%}
\newcommand{\poolMetroGainMaxS}{19}
\newcommand{\poolMetroGainOne}{3.0\%}
\newcommand{\poolMetroGainZero}{5.0\%}
\newcommand{\poolMetroK}{66}
\newcommand{\poolMetroLamLarge}{3.3}
\newcommand{\poolMetroLamSmall}{40}
\newcommand{\poolMetroNaiveOne}{1721.7}
\newcommand{\poolMetroNaiveZero}{1765.4}
\newcommand{\poolMetroPpu}{1349.9}
\newcommand{\poolMetroShadowOne}{1670.7}
\newcommand{\poolMetroShadowZero}{1677.3}
\newcommand{\poolMetroSmax}{40}
\newcommand{\poolMetroSstar}{15.5}
\newcommand{\poolMetroSstarLow}{27.2}
\newcommand{\poolMetroTotHigh}{1421.5}
\newcommand{\poolMetroTotLow}{1109.2}
\newcommand{\poolSAAMax}{0}
\newcommand{\poolSAAMed}{0}
\newcommand{\poolWDROMax}{0}
\newcommand{\poolWDROMed}{0}
\newcommand{\ratioAllHi}{100\%}
\newcommand{\ratioAllLo}{50\%}
\newcommand{\ratioDpThreeHi}{96\%}
\newcommand{\ratioDpThreeLo}{86\%}
\newcommand{\ratioOrcThreeHi}{68\%}
\newcommand{\ratioOrcThreeLo}{52\%}
\newcommand{\reactCui}{21.5}
\newcommand{\reactDet}{179.2}
\newcommand{\reactGounaris}{19.6}
\newcommand{\reactMDRO}{3.6}
\newcommand{\reactSAA}{16.5}
\newcommand{\reactWDRO}{2.9}
\newcommand{\regretAll}{3.8\%}
\newcommand{\regretBoundAll}{7.7\%}
\newcommand{\regretBoundFail}{1}
\newcommand{\regretBoundHolds}{734/735}
\newcommand{\regretCity}{1.9\%}
\newcommand{\regretCleanShare}{90\%}
\newcommand{\regretContained}{100.0\%}
\newcommand{\regretDet}{3.8\%}
\newcommand{\regretFailMax}{0.06\%}
\newcommand{\regretLong}{4.6\%}
\newcommand{\regretMid}{6.8\%}
\newcommand{\regretMin}{1.9\%}
\newcommand{\regretMinRow}{City \textsc{Det}}
\newcommand{\regretRouteMean}{5.1\%}
\newcommand{\regretSaa}{6.8\%}
\newcommand{\regretShort}{2.7\%}
\newcommand{\rlBaton}{33.7\%}
\newcommand{\rlBatonFitS}{1.7}
\newcommand{\rlBatonHo}{32.9\%}
\newcommand{\rlBest}{27.4\%}
\newcommand{\rlInstCI}{[4.4, 7.7]}
\newcommand{\rlInstDelta}{6.0}
\newcommand{\rlInstWin}{12/12}
\newcommand{\rlNInst}{12}
\newcommand{\rlTrainMax}{3}
\newcommand{\rlTrainMin}{1}
\newcommand{\rlWinN}{39/49}
\newcommand{\rlWinP}{$p \le 10^{-3}$}
\newcommand{\salhiActShare}{89\%}
\newcommand{\sbSixtyBaton}{36.8\%}
\newcommand{\sbSixtyHo}{2.1\%}
\newcommand{\sbSixtyThresh}{1.4\%}
\newcommand{\shapeFNine}{0.2\%}
\newcommand{\shapeFSix}{0.2\%}
\newcommand{\shapeFZero}{2.2\%}
\newcommand{\shapePowTenAll}{64\%}
\newcommand{\shapePowTenDayfac}{47\%}
\newcommand{\shapePowTenNine}{88\%}
\newcommand{\shapePowTenSix}{72\%}
\newcommand{\shapePowTenThree}{58\%}
\newcommand{\shapePowTenZero}{51\%}
\newcommand{\shapePowTwentyFiveAll}{89\%}
\newcommand{\shapePowTwentyFiveDayfac}{78\%}
\newcommand{\shapePowTwentyFiveNine}{97\%}
\newcommand{\shapePowTwentyFiveSix}{97\%}
\newcommand{\shapePowTwentyFiveThree}{89\%}
\newcommand{\shapePowTwentyFiveZero}{82\%}
\newcommand{\shapeViolDayfac}{0.6\%}
\newcommand{\shapeViolNine}{0.0\%}
\newcommand{\shapeViolSix}{0.1\%}
\newcommand{\shapeViolThree}{0.1\%}
\newcommand{\shapeViolZero}{0.2\%}
\newcommand{\shopTravelSaving}{25.9\%}
\newcommand{\statCompDeltaMax}{10.3}
\newcommand{\statCompDeltaMin}{0.0}
\newcommand{\statCompLoMin}{-0.4}
\newcommand{\statCompWinMin}{21}
\newcommand{\statConsDeltaMax}{10.3}
\newcommand{\statConsDeltaMin}{6.5}
\newcommand{\statConsPmax}{$p_{\mathrm{Holm}} < 10^{-3}$}
\newcommand{\statConsWinMin}{31}
\newcommand{\statCuiCI}{[+8.1, +12.3]}
\newcommand{\statCuiComp}{two-lever}
\newcommand{\statCuiDelta}{+10.3}
\newcommand{\statCuiDpDelta}{-2.5}
\newcommand{\statCuiP}{0.00}
\newcommand{\statCuiWin}{40}
\newcommand{\statDetCI}{[-0.4, +0.4]}
\newcommand{\statDetComp}{DP$^3_N$}
\newcommand{\statDetDelta}{+0.0}
\newcommand{\statDetDpDelta}{-1.6}
\newcommand{\statDetP}{0.49}
\newcommand{\statDetWin}{21}
\newcommand{\statDpDeltaMax}{-1.6}
\newcommand{\statDpDeltaMin}{-6.6}
\newcommand{\statGounarisCI}{[+7.7, +10.9]}
\newcommand{\statGounarisComp}{two-lever}
\newcommand{\statGounarisDelta}{+9.4}
\newcommand{\statGounarisDpDelta}{-3.1}
\newcommand{\statGounarisP}{0.00}
\newcommand{\statGounarisWin}{40}
\newcommand{\statMDROCI}{[+5.1, +11.1]}
\newcommand{\statMDROComp}{restock}
\newcommand{\statMDRODelta}{+7.9}
\newcommand{\statMDRODpDelta}{-5.0}
\newcommand{\statMDROP}{0.00}
\newcommand{\statMDROWin}{34}
\newcommand{\statPmax}{$p_{\mathrm{Holm}} \le 0.489$}
\newcommand{\statSAACI}{[+8.2, +11.1]}
\newcommand{\statSAAComp}{two-lever}
\newcommand{\statSAADelta}{+9.7}
\newcommand{\statSAADpDelta}{-2.7}
\newcommand{\statSAAP}{0.00}
\newcommand{\statSAAWin}{40}
\newcommand{\statWDROCI}{[+3.6, +9.7]}
\newcommand{\statWDROComp}{restock}
\newcommand{\statWDRODelta}{+6.5}
\newcommand{\statWDRODpDelta}{-6.6}
\newcommand{\statWDROP}{0.00}
\newcommand{\statWDROWin}{31}
\newcommand{\stockoutAll}{47\%}
\newcommand{\stockoutExcess}{10\%}
\newcommand{\svCityBaton}{10.7\%}
\newcommand{\svCityCf}{10.9\%}
\newcommand{\svCityDp}{12.2\%}
\newcommand{\svCityHo}{10.6\%}
\newcommand{\svCityOrc}{36.8\%}
\newcommand{\svCuiBaton}{52.5\%}
\newcommand{\svCuiCf}{53.0\%}
\newcommand{\svCuiDpN}{13.0\%}
\newcommand{\svCuiDpThree}{54.9\%}
\newcommand{\svCuiDpThreeN}{34.3\%}
\newcommand{\svCuiHo}{20.7\%}
\newcommand{\svCuiOrc}{43.1\%}
\newcommand{\svCuiOrcThree}{76.6\%}
\newcommand{\svCuiRoTheta}{18.7\%}
\newcommand{\svCuiThrK}{20.3\%}
\newcommand{\svCuiThresh}{16.7\%}
\newcommand{\svCuiTwo}{42.2\%}
\newcommand{\svDetBaton}{27.9\%}
\newcommand{\svDetCf}{29.2\%}
\newcommand{\svDetDpN}{25.6\%}
\newcommand{\svDetDpThree}{29.5\%}
\newcommand{\svDetDpThreeN}{27.8\%}
\newcommand{\svDetHo}{26.1\%}
\newcommand{\svDetOrc}{40.7\%}
\newcommand{\svDetOrcThree}{47.2\%}
\newcommand{\svDetRoTheta}{24.9\%}
\newcommand{\svDetThrK}{26.0\%}
\newcommand{\svDetThresh}{24.7\%}
\newcommand{\svDetTwo}{26.5\%}
\newcommand{\svGounarisBaton}{53.2\%}
\newcommand{\svGounarisCf}{53.8\%}
\newcommand{\svGounarisDpN}{14.6\%}
\newcommand{\svGounarisDpThree}{56.3\%}
\newcommand{\svGounarisDpThreeN}{36.2\%}
\newcommand{\svGounarisHo}{21.0\%}
\newcommand{\svGounarisOrc}{43.3\%}
\newcommand{\svGounarisOrcThree}{78.0\%}
\newcommand{\svGounarisRoTheta}{19.1\%}
\newcommand{\svGounarisThrK}{20.7\%}
\newcommand{\svGounarisThresh}{16.7\%}
\newcommand{\svGounarisTwo}{43.8\%}
\newcommand{\svMDROBaton}{45.3\%}
\newcommand{\svMDROCf}{45.6\%}
\newcommand{\svMDRODpN}{0.0\%}
\newcommand{\svMDRODpThree}{50.2\%}
\newcommand{\svMDRODpThreeN}{-4.7\%}
\newcommand{\svMDROHo}{11.8\%}
\newcommand{\svMDROOrc}{42.2\%}
\newcommand{\svMDROOrcThree}{79.4\%}
\newcommand{\svMDRORoTheta}{10.3\%}
\newcommand{\svMDROThrK}{8.5\%}
\newcommand{\svMDROThresh}{0.5\%}
\newcommand{\svMDROTwo}{35.6\%}
\newcommand{\svSAABaton}{53.7\%}
\newcommand{\svSAACf}{53.8\%}
\newcommand{\svSAADpN}{12.4\%}
\newcommand{\svSAADpThree}{56.4\%}
\newcommand{\svSAADpThreeN}{35.4\%}
\newcommand{\svSAAHo}{20.1\%}
\newcommand{\svSAAOrc}{43.5\%}
\newcommand{\svSAAOrcThree}{78.4\%}
\newcommand{\svSAARoTheta}{18.4\%}
\newcommand{\svSAAThrK}{19.8\%}
\newcommand{\svSAAThresh}{16.3\%}
\newcommand{\svSAATwo}{43.9\%}
\newcommand{\svSalhiBaton}{53.3\%}
\newcommand{\svSalhiDpThree}{54.3\%}
\newcommand{\svSalhiHo}{42.6\%}
\newcommand{\svSalhiOrc}{49.2\%}
\newcommand{\svWDROBaton}{41.2\%}
\newcommand{\svWDROCf}{41.7\%}
\newcommand{\svWDRODpN}{0.0\%}
\newcommand{\svWDRODpThree}{47.8\%}
\newcommand{\svWDRODpThreeN}{-1.4\%}
\newcommand{\svWDROHo}{10.0\%}
\newcommand{\svWDROOrc}{42.0\%}
\newcommand{\svWDROOrcThree}{79.4\%}
\newcommand{\svWDRORoTheta}{8.9\%}
\newcommand{\svWDROThrK}{7.2\%}
\newcommand{\svWDROThresh}{-1.3\%}
\newcommand{\svWDROTwo}{32.4\%}
\newcommand{\svZpFiftyBaton}{3.5\%}
\newcommand{\svZpFiftyThr}{-1.3\%}
\newcommand{\svZpTwentyFiveBaton}{5.0\%}
\newcommand{\tailBatonHi}{54.3\%}
\newcommand{\tailBatonLo}{22.3\%}
\newcommand{\tailDetBaton}{22.3\%}
\newcommand{\tailDetHo}{23.8\%}
\newcommand{\tailDetThr}{23.6\%}
\newcommand{\thrGapCity}{0.4\%}
\newcommand{\thrGapDet}{1.0\%}
\newcommand{\thrGapLong}{0.9\%}
\newcommand{\thrGapMid}{1.7\%}
\newcommand{\thrGapSaa}{3.0\%}
\newcommand{\thrGapShort}{0.8\%}
\newcommand{\thrkHoGapMax}{3.3}
\newcommand{\thrkHoGapMin}{0.1}
\newcommand{\timeBatonHoMs}{7\,ms}
\newcommand{\timeBatonMs}{24\,ms}
\newcommand{\timeOnlineSkUs}{63\,$\mu$s}
\newcommand{\timeOnlineUs}{1.2\,$\mu$s}
\newcommand{\twoGapHi}{10.3}
\newcommand{\twoGapLo}{1.4}


\setlength{\emergencystretch}{2em}
\definecolor{rev}{RGB}{70,70,70}
\newenvironment{comment}{\par\medskip\noindent\color{rev}\itshape
  \begin{list}{}{\leftmargin=1.2em\rightmargin=1.2em}\item[]}%
  {\end{list}\par}
\newcommand{\resp}{\par\noindent\textbf{Response.}\ }
\newcommand{\chg}{\par\noindent\textbf{Changes.}\ }
\newcommand{\sref}[1]{Section~\ref{m-#1}}
\newcommand{\tref}[1]{Table~\ref{m-#1}}
\newcommand{\fref}[1]{Figure~\ref{m-#1}}
\newcommand{\pref}[1]{Proposition~\ref{m-#1}}
\newcommand{\BATON}{\textsc{Baton}}

\title{Response to the reviewers, second round\\[4pt]
\large Manuscript CAOR-D-26-01885R1, ``Optimal-stopping recourse for
vehicle routing with simultaneous pickup, delivery, and stochastic
demands''}
\author{Quang-Vinh Dang, Hoang-Viet Vu, Phuc-Son Nguyen}
\date{}

\begin{document}
\maketitle

\noindent Dear Editor, dear Reviewers,

We thank the reviewers for their second reading. Reviewer~1 considers
the main concerns resolved. Reviewer~2 asks us to align the
interpretation of \textsc{Baton-cf} and of the fresh-start bias evidence,
and the statistical and implementation descriptions, with what the
evidence shows, and points to two features of the code; Reviewer~3 asks
four questions about the tables and the cost model. We have checked
every implementation point against the code, corrected two conventions
in the evaluation and re-run the affected experiments, and revised the
text where the reviewers found our claims stronger than the evidence.
Section, table and proposition numbers refer to the revised manuscript;
a version with the changes marked against revision~1 is provided
separately.

\section*{Summary of the changes}

\begin{enumerate}[leftmargin=1.4em]
\item \textbf{\textsc{Baton-cf} as an approximation}
  (\sref{sec:actions}). The text now distinguishes the cumulative
  increment $S_k$ from the physical load $X_k$, explains what information
  is retained after unloading, describes \textsc{Baton-cf} as an
  approximation that reduces rather than removes the dependence bias,
  gives a sufficient condition for the monotonicity of $F_k(\cdot)$ and
  checks it empirically.
\item \textbf{Fresh-start bias} (\sref{sec:actions}, \tref{tab:fresh}).
  The three estimands are defined, the two measured differences carry
  instance-level bootstrap intervals, the claim that two biases ``nearly
  cancel'' is withdrawn, and the two reference programs are labelled
  separately.
\item \textbf{Two corrections in the evaluation code}, with the affected
  experiments re-run: routes that depart fully loaded now keep the
  physical slack $B = 0$, and the physical post-return state is capped at
  an empty vehicle. Neither changes a conclusion.
\item \textbf{Scope and statistics}: the day-type experiment is described
  as concerning known day types under exchangeable days, its intervals
  resample instances, ties are reported as the absence of a
  statistically detectable difference, the over-triggering bound is
  stated as a sufficient condition, and the exact dynamic program carries
  a grid-refinement check.
\item \textbf{Reviewer~3's questions}: the meaning of ``bound'' and
  ``reference point'', the contrast between the \textsc{SAA} and
  \textsc{WDRO} gates, relative versus absolute costs in the sensitivity
  analysis (with absolute costs added to \tref{tab:costsens}), and the
  dispatch fee $F_{\mathrm{ho}}$.
\item \textbf{Reproducibility}: the repository state of each submitted
  version is tagged (\texttt{baton-r1}, \texttt{baton-r2}).
\end{enumerate}

% ════════════════════════════════════════════════════════════════════════
\section*{Reviewer 1}

\begin{comment}
I have examined the revised version and consider that the authors have
satisfactorily addressed the main concerns raised in my initial review.
[\ldots] Overall, I consider that the revised manuscript represents a
substantial improvement over the original submission and that my
principal concerns have been adequately resolved.
\end{comment}
\resp We thank the reviewer for the careful second reading and for the
encouraging assessment. The changes made in this round, in response to
Reviewers~2 and~3, narrow several claims further and correct two
conventions in the evaluation code; none of them changes a conclusion
that the reviewer's comments concerned.

% ════════════════════════════════════════════════════════════════════════
\section*{Reviewer 2}

\begin{comment}
1. The remaining information-state approximation in \textsc{Baton-cf}.
[\ldots] Let $S_k$ denote cumulative observed increments without resets
and $X_k$ the physical load after returns; following a return at $r$,
$X_k = S_k - S_r$. Under factor dependence, identical $(k, X_k)$ can
conceal different $S_r$ and different predictions of remaining demand.
[\ldots] The public implementation previously inspected trains
conditional models on unreset cumulative sums but supplies reset $W$
online; if this also applies to the submitted version, \textsc{Baton-cf}
should remain explicitly described as an approximation [\ldots]. The
monotonicity of $F_k$ also needs a condition or separate justification
[\ldots]. Constancy under independence is a population property, not a
guarantee for finite-sample fits or exact policy agreement with
\textsc{Baton}.
\end{comment}
\resp The reviewer's reading of the code is correct, and it applies to
the submitted version: the continuation and fresh-start models of
\textsc{Baton-cf} are fitted on paths indexed by the unreset cumulative
increment, and after a return they are evaluated at the reset (physical)
load. We have rewritten the description in \sref{sec:actions} along the
lines the reviewer suggests. The text now introduces $S_k$ and $X_k$
explicitly and notes that unloading destroys no information: the
operator still observes $S_k$. It explains that under independence the
future depends on $X_k$ alone, whereas under factor dependence the
factor posterior depends on $S_k$, so equal $(k, X_k)$ can hide different
predictions. \textsc{Baton-cf} conditions the price of a \emph{first}
return on its pre-return state ($S_k$), but applies its continuation
models after a return to $X_k$, as \textsc{Baton} does; it is therefore
described as an approximation that reduces bias~(iii) for the first
return without conditioning later decisions on the full history, and an
exact treatment is identified as one on the two-dimensional state
$(k, X_k, S_k)$. All statements that \textsc{Baton-cf} ``removes'' the
dependence bias, or removes it ``at its source'', have been replaced by
``reduces''.

On monotonicity, the reviewer is right that \pref{prop:monotone} does not
cover the conditional estimator. The text now gives a sufficient
condition (the posterior law of the remaining increments stochastically
increasing in $S_k$, and the cost of the remainder of the day under the
downstream policy non-decreasing in those increments, path by path),
notes that the second part need not hold for a fitted policy, and
describes the isotonic form of $\widehat F_k$ as a modelling choice. We
then test it with the binned shape test used for $C_k$, applied to the
realised post-return cost under \textsc{Baton-cf}'s fitted downstream
policy on $5\times10^4$ paths per route: significant decreases occur in
\shapeFZero{} of adjacent-bin pairs under independence, where
$F_k(\cdot)$ is constant in population and the rate is the
false-positive level of the test, and in \shapeFSix{} and \shapeFNine{}
of pairs at $\rho = 0.6$ and $0.9$. Finally, the text now states that
constancy under independence is a population property: in finite
samples the two fits differ, and \textsc{Baton-cf} and \textsc{Baton}
need not take the same decisions.
\chg \sref{sec:actions}, paragraphs on the fresh-start value and on
deployment selection, and after Eq.~\eqref{m-eq:bellman3}; new paragraph
on the monotonicity of the conditional fresh-start value in
\sref{sec:robust}; wording in \sref{sec:setup} and \sref{sec:robust}.

\begin{comment}
2. What Table A.3 establishes about fresh-start bias. [\ldots] the claim
that the first two biases ``nearly cancel'' is not established. For the
same fitted downstream policy, the in-sample and out-of-sample errors are
$+0.50\%$ and $+0.09\%$ of the handoff price, respectively; their
$+0.41\%$ difference does not identify the proposed downward in-sample
effect. I suggest defining the three estimands and reporting uncertainty
for paired differences, or limiting the conclusion [\ldots]. The
population cost of a fixed feasible downstream policy bounds the
optimum, whereas a finite-sample estimate need not; the high-data program
should also be identified as an approximate reference. Finally,
reconcile the references with DP$^3_{50k}$ in Section 3.6: the public
code previously inspected uses constant-F \texttt{fit\_dp\_actions} for
the bias and replacement-F columns, but \texttt{fit\_dp\_actions\_cf} for
the final performance column.
\end{comment}
\resp We agree on every point. The difference of $+0.41\%$ indeed has the
opposite sign to the in-sample optimism we had invoked, so ``nearly
cancel'' was not supported, and we have withdrawn it. Following the
reviewer's suggestion, \sref{sec:actions} now defines three quantities:
the optimal $F_k$; the population value $F^{\mathrm{pop}}_k$ of the reset
state under the fitted downstream policy; and the in-sample estimate
$\widehat F_k$. The text states that $F^{\mathrm{pop}}_k \ge F_k$ is a
population statement that a finite-sample estimate need not respect, and
that the sign of the in-sample error $\widehat F_k - F^{\mathrm{pop}}_k$
is not determined a priori. \tref{tab:fresh} now reports the two
measurable differences with 95\% bootstrap intervals that resample
instances: the in-sample effect $\widehat F_k - F^{\mathrm{pop}}_k$ is
\frSalhiInEff{} of the handoff price on Salhi--Nagy (\frSalhiInEffCI) and
\frDethInEff{} on Dethloff (\frDethInEffCI); the difference
$F^{\mathrm{pop}}_k - F^{\mathrm{ref}}_k$ is \frSalhiDsEff{}
(\frSalhiDsEffCI) and \frDethDsEff{} (\frDethDsEffCI). Because
$F^{\mathrm{ref}}_k$ is itself an estimate, the second difference
combines downstream suboptimality with reference error and identifies
neither, which the text now says. On the city routes this difference
is negative, \frCityDsEff{} (\frCityDsEffCI): the fitted policy does
better than the reference there, which confirms that the reference is
only approximate; the text reports this as well. The conclusion we draw is limited, as
the reviewer proposes, to the small change in policy savings when
$\widehat F_k$ is replaced: on Salhi--Nagy the saving moves from
\frSalhiThree{} to \frSalhiFout{} with $F^{\mathrm{pop}}_k$ and to
\frSalhiFstar{} with $F^{\mathrm{ref}}_k$ (\frSalhiDeltaF{} points,
\frSalhiDeltaFCI).

The reviewer's reading of the code is again correct. The bias and
replacement columns use the \emph{constant-F} high-data program
(\texttt{fit\_dp\_actions}), while the last column, like the main tables,
uses the state-conditional program (\texttt{fit\_dp\_actions\_cf}). The
two are now labelled separately: the reference for the bias columns is
called $F^{\mathrm{ref}}_k$ and is described as the fresh-start value of
DP$^{3,c}_{50\mathrm{k}}$, the unconditional counterpart of
DP$^3_{50\mathrm{k}}$, and both are described as approximate references.
\chg \sref{sec:actions} (paragraph on the estimate of the fresh-start
value); \tref{tab:fresh} (new first two columns, intervals, caption);
paragraph on the fresh-start value in \sref{sec:robust}.

\begin{comment}
3. Scope and uncertainty of the new promotion-day experiment. [\ldots] a
fixed ordinary/promotional mixture can remain exchangeable across days.
I suggest narrowing the abstract and section claims accordingly; this
experiment does not establish robustness to temporal drift or serial
dependence [\ldots]. The interval pooling 80 Det/SAA plans from 40
instances should analyse each gate separately or resample by instance
with both plans retained together. For the deliver-only comparison, the
Table 5 interval $[-6.9, +2.4]$ percentage points supports ``no
statistically detectable difference,'' rather than an established
performance tie.
\end{comment}
\resp We agree. The abstract, the introduction and \sref{sec:robust} now
speak of \emph{known day types} rather than non-exchangeable days, and
the section states that the days are drawn independently from a fixed
mixture, so they remain exchangeable, and that the experiment focuses
on known day types, with temporal drift and serial dependence left for
future work; the conclusion says the same. The interval now resamples instances with both
plans kept together, and we also report each gate separately: the
day-type-specific fits differ from the pooled fit by
\dtAwarePooledDelta{} points overall (\dtAwarePooledCI),
\dtAwarePooledDetDelta{} on \textsc{Det} plans (\dtAwarePooledDetCI) and
\dtAwarePooledSaaDelta{} on \textsc{SAA} plans (\dtAwarePooledSaaCI).
Finally, wherever an interval contains zero we now write that the
difference is not statistically detectable rather than that the
policies are tied, including the deliver-only comparison, the
deterministic gate and the dependence rows.
\chg Abstract; introduction; \sref{sec:robust} (``Known day types'');
conclusion; caption of \tref{tab:daytype}; wording throughout
\sref{sec:study}.

\begin{comment}
4. Physical-reset conventions and evaluation code. [\ldots] I suggest
explaining what a known departure load reveals about future deliveries
when actual quantities are loaded, or how deliveries exceeding loaded
cargo are handled when nominal quantities are used, so that
$L_0 - D_{\le k}$ has a consistent physical meaning. The public
\texttt{run\_realistic\_eval.py} and \texttt{route\_setup} in
\texttt{run\_baton\_r1.py} previously inspected replace $B = Q - L_0$
with an empirical peak quantile when $B \le 0$, although zero slack is a
valid fully loaded departure. State whether this branch is activated in
the reported experiments and how many routes it affects. [\ldots] A
fixed R1 repository commit or archive would allow verification against
the submitted implementation.
\end{comment}
\resp Thank you; both points led to corrections. On the loading
convention, our simulator uses \emph{nominal} loading: the vehicle
departs with the planned delivery volume $L_0 = \sum_k \bar d_k$ and the
realised deliveries are revealed on service. \sref{sec:setting} now
states this, reports how often the day's deliveries exceed the cargo
loaded (on \stockoutAll{} of route-days, by \stockoutExcess{} of $L_0$ on
average), and explains that such shortfalls are outside the model:
they are not priced and no policy affects them, so they do not change
the comparison of policies. It also explains that loading the realised
volume would make $L_0$ random but known at departure and reveal the
day's delivery total before the first stop, a different information
structure that we do not study. For the physical reset this means that
the undelivered cargo is $\max\{L_0 - D_{\le k}, 0\}$; the submitted
implementation did not cap it, so the physical post-return state could
fall below $-L_0$. The state is now $x_k = -\min\{D_{\le k}, L_0\}$, and
the physical-reset column of \tref{tab:fresh} was re-computed with it
(\frSalhiExact{} on Salhi--Nagy and \frDethExact{} on Dethloff). The
conclusion that the conservative reset understates the value of the
return is unchanged.

On the slack, the branch was activated. No route of any reported
experiment departs over capacity, but 16 routes depart exactly full
($B = 0$): one Salhi--Nagy route and fifteen routes of the city
instances and their twins (none of the Dethloff routes, none of the
routes of the RL comparison, and none of the city \textsc{SAA} plans).
For these routes the code replaced $B = 0$ by an empirical peak
quantile. We now keep the physical slack $B = 0$ in all evaluation
code, so that such a route breaches as soon as its running load rises,
and we re-ran every affected experiment: the Salhi--Nagy and city
benchmarks of \tref{tab:large}, the fresh-start, over-triggering,
data-budget and standby-pool experiments. The mean savings of the
implementable policies moved by at most 0.7 percentage points (on the
uniform city twins) and those of the reference points by at most 1.7
points (the three-action clairvoyant on the 50\% deliver-only twins); no
ranking or conclusion changed, and the manuscript reports the re-run
numbers. Last,
the repository now carries the tag \texttt{baton-r1} for the state that
produced revision~1, and \texttt{baton-r2} for this revision, so that the
code can be checked against each submitted version; the Data
availability statement says so.
\chg \sref{sec:setting} (loading convention, fully loaded departures);
\sref{sec:actions} (physical reset); \tref{tab:large}, \tref{tab:fresh},
\tref{tab:regret}, \tref{tab:pool} and \tref{tab:budget} re-computed;
Data availability.

\begin{comment}
(a) Proposition 3 and Table 6 address over-triggering. I suggest
replacing ``good enough exactly when'' with a sufficient-condition
statement: a small upper bound guarantees a small loss but is not
necessary. Interpret the numerical verification with the qualification
that the policy labelled optimal is estimated from 50,000 paths.
\end{comment}
\resp Done. The text after \pref{prop:regret} now presents the bound as a
sufficient condition for the myopic threshold to be good enough, and
says that it is not necessary because the bound need not be tight. The
discussion of \tref{tab:regret} and its caption now state that the rule
labelled optimal is itself an estimate from $5\times10^4$ paths, so that
the table verifies the proposition up to that estimation error.
\chg Paragraph after \pref{prop:regret}; \sref{sec:learning};
caption of \tref{tab:regret}.

\begin{comment}
(b) Table 8 is welcome. Since exact-optimum percentages now appear in
the abstract, I suggest adding a small grid-refinement or truncation
check for the optimal values. Agreement between simulated and grid-based
reactive costs alone does not establish the error in the optimal policy
value.
\end{comment}
\resp We agree, and we have added the check. For all \gridRoutes{}
routes of \tref{tab:exact} we solved the handoff-only and the
three-action programs on grids of 1{,}500, 3{,}000 and 6{,}000 load cells,
and with the lower truncation and the support of the demand
distributions widened (from 6 to 10 standard deviations and from 12 to
20 coefficients of variation). Values of single routes move by up to \gridCoarseThreeMax{} between
1{,}500 and 6{,}000 cells (median \gridCoarseThreeMedian; \gridMidThreeMax{}
between 3{,}000 and 6{,}000 cells), but the discretisation affects the
reactive, handoff-only and three-action values alike, so the pooled
exact savings of \tref{tab:exact}, and with them the percentages in the
abstract, change by at most \gridPooledMaxDiff{} percentage points across
all grids and truncations. The text and the caption of \tref{tab:exact}
now report this check.
\chg Paragraph on the exact benchmark in \sref{sec:robust}; caption of
\tref{tab:exact}.

\begin{comment}
(c) I suggest stating Algorithm 1's action-specific updates explicitly,
distinguishing realised continuation-path costs from the estimated
$R_k + F_k$ inserted upon return. Also acknowledge that the equal-data
three-action DP splits its training sample while Baton reuses it, so the
comparisons reflect sample allocation as well as estimator differences.
\end{comment}
\resp Done. Algorithm~\ref{m-alg:offline} now spells out the three
updates: a path whose estimated continuation cost exceeds the handoff
price is assigned that price; a path for which the return is cheapest is
assigned the estimated $R_k + F_k$; and a path that continues keeps its
realised cost along the path. The protocol paragraph in \sref{sec:setup}
now notes that DP$^3_N$ splits its $10^3$ training days between its two
estimation steps whereas \textsc{Baton} uses all of them for both, so
that the comparison reflects the allocation of the sample as well as the
estimators.
\chg Algorithm~\ref{m-alg:offline}; ``Protocol'' paragraph in
\sref{sec:setup}.

% ════════════════════════════════════════════════════════════════════════
\section*{Reviewer 3}

\begin{comment}
1. In Table 1, both oracle-ho and oracle3 are labeled as `bounds',
however, in subsequent simulation studies, as shown in Table 2, Table 11,
Fig. 5, and Fig. 8, only oracle3 appears as a bound, oracle-ho becomes a
reference point, is the reference point equivalent to bound?
\end{comment}
\resp Thank you for pointing out the inconsistency. The two words are
not equivalent, and we now use them consistently. A clairvoyant policy
is a \emph{bound} for a class of policies when no policy in that class
can do better on any day. The handoff-only clairvoyant (oracle-ho) knows
the day's demands but may only hand off; it is therefore a bound for the
handoff-only policies (the thresholds, the rule-based triggers, the
rollouts, \textsc{Baton-ho}), but not for policies that may also return
to the depot, which can prevent a breach more cheaply than any handoff
and can therefore exceed it, as \textsc{Baton} does on four of the six
gates. For those policies oracle-ho is only a \emph{reference point}.
The three-action clairvoyant (oracle$^3$) is a bound for every policy in
the paper. \tref{tab:competitors} now says for which policies each
clairvoyant is a bound, \sref{sec:algorithms} defines the two terms, and
the captions of the results tables use them in this sense; the
reference-point columns of the tables group the clairvoyant policies
with the high-data programs, which are neither bounds nor competitors.
\chg \tref{tab:competitors}; \sref{sec:algorithms}; caption of
\tref{tab:grand}.

\begin{comment}
2. In Sec 4.1 planning regimes, it claims that `the WDRO gate is the SAA
gate at a tighter threshold', in Table 2, it can be seen that in SAA
gate, the value of BATON (53.7) is much higher than that of oracle-ho
(43.5), while in WDRO gate, the value of BATON (41.2) is almost the same
as that of oracle-ho (42.0). Meanwhile, the value of BATON in SAA gate
(53.7) is much higher than that of WDRO gate (41.2). What are the
explanation for these?
\end{comment}
\resp The values in Table~2 are savings in per cent of each gate's own
reactive cost, and that cost differs by a factor of about five between
the two gates. The tighter threshold of the \textsc{WDRO} gate produces
much safer plans: the reactive policy costs \reactWDRO{} per plan-day on
\textsc{WDRO} plans against \reactSAA{} on \textsc{SAA} plans, and at
least one route of a plan breaches on \pemWDRO{} of days against
\pemSAA. In absolute terms \BATON{} is far cheaper on \textsc{WDRO} plans
(\costWDROBaton{} against \costSAABaton{} per plan-day), but there are
fewer breaches to prevent, and a thousand training days contain few of
them; this is the rare-breach regime in which \BATON{} falls furthest
short of what more data would allow (\svWDROBaton{} against
\svWDRODpThree{} for the high-data three-action program, compared with
\svSAABaton{} against \svSAADpThree{} on \textsc{SAA}). The handoff-only
clairvoyant knows the day's demands and saves a similar share on both
gates (\svWDROOrc{} and \svSAAOrc), so \BATON{} exceeds it by a wide
margin on \textsc{SAA} plans but not on \textsc{WDRO} plans. We have
added this explanation to \sref{sec:headline}.
\chg New paragraph in \sref{sec:headline}.

\begin{comment}
3. In Table 11, it shows that the value of BATON with a cheap standby
(Fsb=10, BATON=52.7) is much higher than that with a dear standby
(Fsb=35, BATON=42.6), while the value of BATON with a cheap emergencies
(Femg=25, BATON=40.1) is much low than that with a dear emergencies
(Femg=60, BATON=55.2). Does this mean that cheap standby is better, and
cheap emergencies are not better?
\end{comment}
\resp No: in absolute terms both are better. The values in
\tref{tab:costsens} are savings relative to the reactive cost of the
same configuration, and the reactive cost itself depends on the
emergency price. \BATON's expected daily cost per plan is \csBaseBaton{}
at the default prices, \csSbLowBaton{} with cheap standby and
\csEmgLowBaton{} with cheap emergencies, against \csSbHighBaton{} with
dear standby and \csEmgHighBaton{} with dear emergencies. The savings
move in opposite directions for a different reason. Cheap standby makes
the interventions cheaper while the reactive cost stays at
\csBaseReact{}, so the saving rises. Cheap emergencies lower the
reactive cost itself, to \csEmgLowReact, so there is less to save, and
dear emergencies raise it, to \csEmgHighReact, so intervention recovers
a larger share of a larger bill. We have added the reactive and
\BATON{} daily costs as two columns of \tref{tab:costsens} and this
explanation to \sref{sec:sensitivity}.
\chg \tref{tab:costsens} (two new columns, caption);
\sref{sec:sensitivity}.

\begin{comment}
4. In Eq. (4), does the Fho represent dispatch fee?
\end{comment}
\resp Yes. $F_{\mathrm{ho}}$ is the dispatch fee, the fixed charge for
calling out a standby vehicle. The symbol was not named in the text,
which now says so where Eq.~\eqref{m-eq:handoffprice} is introduced.
\chg \sref{sec:setting}.

\bigskip
We thank the reviewers again for their time and for comments that have
made the paper more precise.

\medskip
\noindent Sincerely,\\
Quang-Vinh Dang (corresponding author), on behalf of all authors

\end{document}
